\documentclass[11pt]{letter}
\usepackage[margin=1in]{geometry}
\usepackage{hyperref}

\signature{
\c{C}a\u{g}r{\i} \c{S}ahin\\
Department of Software Engineering\\
F{\i}rat University, Elaz{\i}\u{g}, Turkey\\
Email: csahin@firat.edu.tr
}

\address{
\c{C}a\u{g}r{\i} \c{S}ahin\\
Department of Software Engineering\\
F{\i}rat University\\
23119 Elaz{\i}\u{g}, Turkey
}

\begin{document}

\begin{letter}{
Editor-in-Chief\\
IEEE Access\\
IEEE Publishing Operations
}

\opening{Dear Editor,}

We are pleased to submit our manuscript titled ``A Comparative Study of Convolutional and Transformer Architectures Under Adversarial Perturbations: Gradient Characteristics and Transferability Analysis'' for consideration as a Regular Article in \textit{IEEE Access}.

This paper investigates the mechanistic differences underlying adversarial vulnerability in convolutional neural networks (CNNs) and Vision Transformers (ViTs). Rather than merely comparing robustness scores, we ask \textit{why} these architectures behave differently under adversarial perturbations, an increasingly important question as both architecture families are deployed in safety-critical systems. Using CIFAR-10 with adversarially trained ResNet-18 and ViT-Tiny models evaluated under AutoAttack ($\epsilon = 8/255$), we present four key findings:

\begin{itemize}
    \item \textbf{Transfer attack asymmetry:} CNN-generated adversarial examples transfer to ViTs at a 47.5\% success rate, while ViT-to-CNN transfers succeed at only 33.5\%, revealing a 14-percentage-point directional gap with practical implications for black-box attack strategies.
    \item \textbf{Gradient characteristics:} ViT gradients are 2.2$\times$ more aligned across samples and CNN gradients are 4.5$\times$ more sparse, providing a quantitative explanation for the observed transfer asymmetry.
    \item \textbf{Feature degradation analysis:} Layer-wise examination of ViT internals shows progressive semantic degradation under adversarial perturbations, with cosine similarity dropping from 0.997 to 0.917 and 7.86\% feature norm reduction in late layers.
    \item \textbf{Design principles:} Based on our mechanistic findings, we propose specific recommendations for robust hybrid system design.
\end{itemize}

We believe this work is well suited for \textit{IEEE Access} because of the journal's interdisciplinary readership spanning computer vision, machine learning security, and deep learning systems. Our analysis bridges adversarial robustness and architectural understanding, topics that are highly relevant to the diverse IEEE community. The open-access nature of \textit{IEEE Access} will also maximize the impact of these findings for researchers and practitioners working on secure AI deployment.

\textbf{Declarations:}
\begin{itemize}
    \item This manuscript is original work and has not been published elsewhere.
    \item The manuscript is not under consideration by any other journal.
    \item All authors have reviewed and approved the final manuscript.
    \item All experiments use the publicly available CIFAR-10 dataset and standard evaluation protocols.
    \item The authors declare no conflict of interest.
    \item Code and trained models will be made publicly available upon acceptance.
\end{itemize}

Thank you for considering our manuscript. We look forward to your editorial decision.

\closing{Sincerely,}

\end{letter}
\end{document}
